
\begin{frame}{L'objectif}

Nous allons travailler sur des systèmes de gestion de données \alert{distribués}. 

\vfill

Un système distribué fonctionne, de manière coordonnée, sur plusieurs machines, ou \alert{serveurs}.

\vfill

Dans ce préliminaire au cours, nous allons étudier Docker, un outil qui permet de simuler un
système distribué, avec très peu de ressources, sur un seule ordinateur (par exemple votre portable).

\vfill

\begin{beamerboxesrounded}{Votre objectif: maîtriser l'installation d'un système avec Docker}
Vous devez pouvoir installer en quelques clics un système avec Docker, et comprendre
son architecture.
\end{beamerboxesrounded}

\end{frame}


\subsection{Rappels}



\begin{frame}{Système distribué: rappels}

\begin{redbullet}
  \item  un \alert{serveur machine} est un ordinateur, tournant sous un système d'exploitation,
    et connecté en permanence au réseau via des ports; il a une IP (adresse internet)
  \item   un \alert{serveur logiciel} est un processus exécuté en tâche de fond d'un serveur machine qui
    communique avec des \alert{clients (logiciels)} via un port particulier.
  \item Une \alert{machine virtuelle} est un programme qui simule, sur une machine hôte,
     un autre ordinateur.
   \item Un \alert{client} (logiciel) est un programme qui communique avec un serveur (logiciel)
\end{redbullet}
    
\begin{defblock}{Exemple}
Un serveur web est un processus (Apache par exemple) qui communique sur le
port 80 d'un serveur machine. Si ce serveur machine a pour IP 163.12.9.10, alors tout
client web peut s'adresser au serveur web à l'adresse 163.12.9.10:80.
\end{defblock}

\end{frame}


\begin{frame}{Illustration}

Voici un exemple de système distribué.

\vfill

\figSlideWithSize{Serveurs}{10cm}

\end{frame}

\begin{frame}{Docker}

Docker permet de créer des \alert{conteneurs} Linux, des machines virtuelles extrêmement légères.
\vfill

\figSlideWithSize{schema-docker}{9cm}

\end{frame}


\begin{frame}{Vocabulaire}

\begin{redbullet}
  \item Le \alert{système hôte} est le système d'exploitation principal gérant votre machine ;
  \item  \alert{Machine Docker} (ou \alert{Docker VM}), désigne le système Linux où Docker s'exécute; 
  \item \alert{Docker engine} est le programme qui gère les conteneurs; il s'exécute
    dans la machine Docker.
  \item Un \alert{conteneur} est une partie autonome de la \alert{Docker VM}, se comportant comme une
    machine indépendante.
  \item Le \alert{client Docker} est l'utilitaire  grâce auquel on transmet au moteur 
    les commandes de gestion de ces conteneurs. 
 \end{redbullet}

 Le client Docker peut être soit la ligne de commande (\alert{Docker CLI}), soit
 une interface comme \alert{Kitematic}.

\end{frame}


\begin{frame}{Les images}

Image = système pré-configuré, prêt à l'emploi pour un service particulier (par exemple un serveur
MySQL, ou MongoDB).

\vfill

\figSlideWithSize{images-docker}{8cm}

\vfill
Une image s'installe très facilement, et se place dans un conteneur.

\end{frame}


\begin{frame}{Comment communiquer avec un conteneur ?}

J'ai lancé un conteneur MySQL. Mon serveur tourne, sur le port 3306, \alert{dans le conteneur}.

\vfill

\alert{Question:} comment mon client MySQL peut-il communiquer avec ce serveur?

\vfill

\alert{Réponse: il faut associer un port de la machine Docker avec le port 3306 du conteneur} (\textit{port forwarding}).

\vfill

Par exemple, 
\begin{redbullet}
 \item je vais associer  le port 6603 de la machine Docker avec le port 3306 de \texttt{serveur-mysql}
 \item je vais associer le port 6604 de la machine Docker avec le port 3306 de \texttt{autre-serveur}
\end{redbullet}


\end{frame}

\begin{frame}{À l'action !}

Et voilà pour les principes.

\vfill

Pour la suite : démonstration de l'installation d'un système.

\vfill

\alert{Votre mission}: effectuer les manipulations Docker décrites dans le support de cours, jusqu'à ce que vous 
soyez à l'aise \alert{en comprenant bien ce qui se passe.}


\end{frame}
